% =============================================================
%  FPGA与VHDL考试——章节变式题详解
%  覆盖主教材每个章节中所有"考试常见变式"的完整教学
%  严格遵循：电路→原理→状态变化→设计思想→模板→VHDL→变式
% =============================================================

\documentclass[12pt,a4paper,openany]{ctexbook}

\usepackage{fontspec}
\setmainfont{Times New Roman}
\setCJKmainfont{SimSun}[AutoFakeBold=2.5]
\setCJKsansfont{SimHei}

\usepackage[top=2.5cm,bottom=2.5cm,left=2.5cm,right=2.5cm]{geometry}
\usepackage{amsmath,amssymb}
\usepackage{tikz,circuitikz}
\usetikzlibrary{arrows,shapes,positioning,calc,automata,patterns}

\usepackage{tabularx,booktabs,multirow,makecell,colortbl}
\usepackage{listings,xcolor,enumitem,caption,float,framed}

% Simple VHDL listings (no custom language to avoid TeX conflicts)
\definecolor{vhdlbg}{rgb}{0.97,0.97,0.97}
\lstset{language={},basicstyle=\ttfamily\footnotesize,
  numbers=left,numberstyle=\tiny\color{gray},
  backgroundcolor=\color{vhdlbg},frame=single,tabsize=2,
  breaklines=true,showstringspaces=false}

\usepackage[colorlinks=true,linkcolor=blue!60!black,urlcolor=blue!60!black]{hyperref}

% Simple framed boxes
\newenvironment{detailbox}[1][详解]
  {\begin{oframed}\textbf{\large #1}\\[3pt]}
  {\end{oframed}}

\newenvironment{examfocus}[1][考试聚焦]
  {\begin{oframed}\textbf{\textcolor{red!70!black}{#1}}\\[3pt]}
  {\end{oframed}}

\title{\Huge\textbf{FPGA与VHDL考试\\章节变式题详解}\\[0.5cm]
       \Large 覆盖主教材全部23个电路的\\所有考试变式}
\author{}
\date{\today}

\begin{document}
\maketitle
\tableofcontents

% ============================================================
\part{组合逻辑电路变式详解}
\chapter{8-3编码器——5大变式详解}

\section{变式1：16-4编码器}

\begin{detailbox}[完整教学]
\textbf{【电路升级】}
8-3编码器：$2^3=8$个输入→3位输出。
16-4编码器：$2^4=16$个输入→4位输出。

\textbf{【设计思想】}
原理完全不变——每个输出位=所有编号中该位为1的输入的OR。

对于16-4编码器，输出$Y_3Y_2Y_1Y_0$，输入$I_0\sim I_{15}$。

\textbf{【逻辑推导】}
\begin{align}
Y_3 &= I_8+I_9+I_{10}+I_{11}+I_{12}+I_{13}+I_{14}+I_{15} \quad\text{(编号8-15的bit3=1)}\\
Y_2 &= I_4+I_5+I_6+I_7+I_{12}+I_{13}+I_{14}+I_{15} \quad\text{(bit2=1的编号)}\\
Y_1 &= I_2+I_3+I_6+I_7+I_{10}+I_{11}+I_{14}+I_{15}\\
Y_0 &= I_1+I_3+I_5+I_7+I_9+I_{11}+I_{13}+I_{15}
\end{align}

\textbf{【通用公式】}对于$2^n$-$n$编码器：
$$Y_k = \sum_{i=0}^{2^n-1} I_i \quad\text{当} i \text{的第}k\text{位为1时}$$

\textbf{【VHDL】}
\begin{lstlisting}
entity encoder_16to4 is
  port (I : in std_logic_vector(15 downto 0);
        Y : out std_logic_vector(3 downto 0));
end encoder_16to4;

architecture behavioral of encoder_16to4 is
begin
  process(I)
  begin
    case I is
      when "0000000000000001" => Y <= "0000";
      when "0000000000000010" => Y <= "0001";
      -- ... 中间省略 ...
      when "1000000000000000" => Y <= "1111";
      when others => Y <= "0000";
    end case;
  end process;
end behavioral;
\end{lstlisting}

\textbf{【更好的写法——用for循环查找】}
\begin{lstlisting}
process(I)
begin
  Y <= "0000";
  for i in 0 to 15 loop
    if I(i) = '1' then
      Y <= std_logic_vector(to_unsigned(i, 4));
    end if;
  end loop;
end process;
\end{lstlisting}

\textbf{【考试聚焦】}$2^n$-$n$编码器的通用公式。16-4编码器需要16个case分支。
\end{detailbox}


\section{变式2：优先级编码器（必考！）}

\begin{detailbox}[完整教学]

\textbf{【问题】}
普通编码器假设同一时刻只有一个输入为1。但当多个输入同时为1时（比如多个按键同时按下），普通编码器产生错误输出。需要\textbf{优先级}——编号大的输入优先。

\textbf{【设计思想】}
用if-else级联：从高优先级到低优先级依次判断。I7优先级最高，I0最低。

\textbf{【逻辑推导】}
\begin{center}
\begin{tabular}{|c|c|c|c|c|c|c|c|c||c|c|c|}
\hline
I7 & I6 & I5 & I4 & I3 & I2 & I1 & I0 & Y2 & Y1 & Y0 \\
\hline
1 & x & x & x & x & x & x & x & 1 & 1 & 1 \\
0 & 1 & x & x & x & x & x & x & 1 & 1 & 0 \\
0 & 0 & 1 & x & x & x & x & x & 1 & 0 & 1 \\
0 & 0 & 0 & 1 & x & x & x & x & 1 & 0 & 0 \\
0 & 0 & 0 & 0 & 1 & x & x & x & 0 & 1 & 1 \\
0 & 0 & 0 & 0 & 0 & 1 & x & x & 0 & 1 & 0 \\
0 & 0 & 0 & 0 & 0 & 0 & 1 & x & 0 & 0 & 1 \\
0 & 0 & 0 & 0 & 0 & 0 & 0 & 1 & 0 & 0 & 0 \\
0 & 0 & 0 & 0 & 0 & 0 & 0 & 0 & 0 & 0 & 0 \\
\hline
\end{tabular}
\end{center}
（x=任意值，不管它是0还是1）

\textbf{【VHDL——if-else级联】}
\begin{lstlisting}
entity priority_encoder is
  port (I : in std_logic_vector(7 downto 0);
        Y : out std_logic_vector(2 downto 0);
        GS : out std_logic);  -- Group Select: 有输入=1
end priority_encoder;

architecture behavioral of priority_encoder is
begin
  process(I)
  begin
    if I(7) = '1' then
      Y <= "111"; GS <= '1';
    elsif I(6) = '1' then
      Y <= "110"; GS <= '1';
    elsif I(5) = '1' then
      Y <= "101"; GS <= '1';
    elsif I(4) = '1' then
      Y <= "100"; GS <= '1';
    elsif I(3) = '1' then
      Y <= "011"; GS <= '1';
    elsif I(2) = '1' then
      Y <= "010"; GS <= '1';
    elsif I(1) = '1' then
      Y <= "001"; GS <= '1';
    elsif I(0) = '1' then
      Y <= "000"; GS <= '1';
    else
      Y <= "000"; GS <= '0';  -- 无有效输入
    end if;
  end process;
end behavioral;
\end{lstlisting}

\textbf{【为什么if-else实现优先级】}
if-else的语义：从上到下依次判断，第一个为真的分支被执行。
所以I7在第一个if → I7优先级最高。

\textbf{【综合成什么硬件】}
级联MUX链（优先级MUX树）。I7→I6→...→I0逐级选择。

\textbf{【GS信号的必要性】}
没有GS时，输入全0（无有效输入）和只有I0=1时，输出都是000——无法区分！
GS=1表示"有有效输入"。

\textbf{【Testbench——测试优先级】}
\begin{lstlisting}
-- 同时置I7和I0=1 → 应该输出111(I7优先)
I <= "10000001"; wait for 10 ns;
-- 输出应为111, GS=1 (I7优先级高于I0)
\end{lstlisting}

\begin{examfocus}
\textbf{必考！}优先级编码器是考试最高频的组合逻辑变式之一。
\begin{itemize}
  \item if-else级联 = 优先级递减的硬件
  \item GS信号区分"无输入"和"选择I0"
  \item 优先级从高到低排列if条件
\end{itemize}
\end{examfocus}

\end{detailbox}


\section{变式3：带使能端的编码器}

\begin{detailbox}[完整教学]

\textbf{【电路】}
在标准编码器基础上增加Enable(E)输入。E=1时编码器正常工作；E=0时所有输出为0（或高阻态）。

\textbf{【VHDL】}
\begin{lstlisting}
entity encoder_with_enable is
  port (I : in std_logic_vector(7 downto 0);
        E : in std_logic;
        Y : out std_logic_vector(2 downto 0));
end encoder_with_enable;

architecture behavioral of encoder_with_enable is
begin
  process(I, E)
  begin
    if E = '0' then
      Y <= "000";  -- 禁止输出
    else
      -- 正常编码逻辑
      case I is
        when "00000001" => Y <= "000";
        when "00000010" => Y <= "001";
        -- ...
        when others => Y <= "000";
      end case;
    end if;
  end process;
end behavioral;
\end{lstlisting}

\textbf{【设计思想】}
使能端是数字电路中常见的控制信号。在不使用时关断输出以节省功耗/避免总线冲突。

\end{detailbox}


\section{变式4：带GS（Group Select）输出的编码器}

\begin{detailbox}[完整教学]

\textbf{【问题】}
普通8-3编码器：当所有输入为0时，输出000。但当$I_0=1$时，输出也是000。如何区分这两种情况？

\textbf{【方案】}增加GS（Group Select/Valid）输出：
\begin{itemize}
  \item 任意输入=1 → GS=1（表示输出有效）
  \item 所有输入=0 → GS=0（表示输出无效/无输入选中）
\end{itemize}

\textbf{【VHDL】}
\begin{lstlisting}
GS <= '0' when I = "00000000" else '1';
-- 或
GS <= I(0) or I(1) or I(2) or I(3) or I(4) or I(5) or I(6) or I(7);
\end{lstlisting}

\textbf{【GS硬件】}一个8输入OR门——需要8个输入引脚。

\textbf{【考试常考】}GS信号的意义：解决编码器输出的二义性问题。

\end{detailbox}


\section{变式5：低电平有效的编码器}

\begin{detailbox}[完整教学]

\textbf{【电路变化】}
输入和输出都是\textbf{低电平有效}（0表示有效/选中，1表示无效）。
这在实际IC（如74LS148）中很常见——低电平有效更节省功耗，且与TTL电平兼容。

\textbf{【逻辑推导】}
低电平有效时：
\begin{itemize}
  \item 输入端：$\overline{I_0}, \overline{I_1}, ..., \overline{I_7}$（上划线表示低电平有效）
  \item 当$\overline{I_k}=0$且其他为1时 → 输出编码为k
  \item 用NAND门替代OR门来实现（因为低有效的或=高有效的与非）
\end{itemize}

\textbf{【VHDL】}
\begin{lstlisting}
entity encoder_active_low is
  port (
    I_n : in  std_logic_vector(7 downto 0);  -- 低有效
    Y_n : out std_logic_vector(2 downto 0);  -- 低有效
    GS_n: out std_logic  -- 低有效：0=有输入, 1=无输入
  );
end encoder_active_low;

architecture behavioral of encoder_active_low is
begin
  process(I_n)
  begin
    GS_n <= '1';  -- 默认：无有效输入
    Y_n <= "111";

    if I_n(7) = '0' then
      Y_n <= "000"; GS_n <= '0';  -- 7的编码=111, 反相=000
    elsif I_n(6) = '0' then
      Y_n <= "001"; GS_n <= '0';  -- 6=110, 反相=001
    elsif I_n(5) = '0' then
      Y_n <= "010"; GS_n <= '0';
    elsif I_n(4) = '0' then
      Y_n <= "011"; GS_n <= '0';
    elsif I_n(3) = '0' then
      Y_n <= "100"; GS_n <= '0';
    elsif I_n(2) = '0' then
      Y_n <= "101"; GS_n <= '0';
    elsif I_n(1) = '0' then
      Y_n <= "110"; GS_n <= '0';
    elsif I_n(0) = '0' then
      Y_n <= "111"; GS_n <= '0';  -- 0的编码=000, 反相=111
    end if;
  end process;
end behavioral;
\end{lstlisting}

\textbf{【关键理解】}
低有效编码：输入0有效→输出也是0=111(即7的二进制反码)。高有效和低有效编码器的输出值互为反码。

\textbf{【考试提示】}看清题目要求：高有效还是低有效？74LS148是低有效优先级编码器，这是常考器件。

\end{detailbox}


\section{编码器变式总结}

\begin{examfocus}[编码器5大变式速查]
\begin{center}
\begin{tabular}{|l|l|l|}
\hline
\textbf{变式} & \textbf{核心改动} & \textbf{VHDL关键} \\
\hline
16-4编码 & 规模扩大 & case 16分支/for循环 \\
优先级编码 & 允许同时多输入 & if-else级联(高→低) \\
带使能 & 增加E控制 & if E='0' then Y<="000" \\
带GS输出 & 增加有效标志 & GS<=not(I=全0) \\
低电平有效 & 0有效1无效 & 判断I(n)='0', 输出取反 \\
\hline
\end{tabular}
\end{center}
\end{examfocus}

\chapter{3-8译码器——5大变式详解}

\section{变式1：74LS138风格译码器（低有效输出+三使能端）——必考！}

\begin{detailbox}[完整教学]

\textbf{【电路】}
74LS138是考试中最常出现的译码器芯片。特点：
\begin{itemize}
  \item 3位地址输入A2,A1,A0
  \item 8个输出Y0-Y7，\textbf{全部低电平有效}（选中=0，未选中=1）
  \item 3个使能端：G1(高有效), G2A(低有效), G2B(低有效)
  \item 只有三个使能端同时有效（G1=1, G2A=0, G2B=0）时，译码器才工作
\end{itemize}

\textbf{【使能逻辑】}$enable = G1 \cdot \overline{G2A} \cdot \overline{G2B}$

\textbf{【VHDL】}
\begin{lstlisting}
entity decoder_74ls138 is
  port (
    A    : in  std_logic_vector(2 downto 0);
    G1   : in  std_logic;
    G2A_n: in  std_logic;  -- 低有效
    G2B_n: in  std_logic;  -- 低有效
    Y_n  : out std_logic_vector(7 downto 0)  -- 低有效
  );
end decoder_74ls138;

architecture behavioral of decoder_74ls138 is
  signal enable : std_logic;
begin
  enable <= G1 and (not G2A_n) and (not G2B_n);

  process(A, enable)
  begin
    if enable = '0' then
      Y_n <= "11111111";  -- 全高=全无效
    else
      case A is
        when "000" => Y_n <= "11111110";  -- Y0=0
        when "001" => Y_n <= "11111101";  -- Y1=0
        when "010" => Y_n <= "11111011";
        when "011" => Y_n <= "11110111";
        when "100" => Y_n <= "11101111";
        when "101" => Y_n <= "11011111";
        when "110" => Y_n <= "10111111";
        when "111" => Y_n <= "01111111";
        when others => Y_n <= "11111111";
      end case;
    end if;
  end process;
end behavioral;
\end{lstlisting}

\textbf{【为什么用3个使能端】}
多个使能端方便级联。G2A和G2B是低有效，G1是高有效——这种混合设计使得可以用不同信号控制，增加灵活性。

\textbf{【常考陷阱】}
\begin{itemize}
  \item 看清输出是低有效还是高有效——74LS138输出0表示选中
  \item 三个使能端必须\textbf{同时}满足条件译码器才工作
  \item 使能无效时所有输出为1（全高=全无效）
\end{itemize}

\end{detailbox}


\section{变式2：用译码器实现任意逻辑函数（必考！）}

\begin{detailbox}[完整教学]

\textbf{【核心原理】}
译码器的每个输出=一个最小项。\textbf{任何组合逻辑函数}=若干最小项之和（OR）。

\textbf{【通用公式】}$F(A,B,C) = \sum m(\dots)= Y_{i1}+Y_{i2}+Y_{i3}+\dots$

\textbf{【例题】}用3-8译码器+一个OR门实现$F(A,B,C)=\sum m(1,3,5,7)$。

\textbf{【解法】}
\begin{enumerate}
  \item 译码器产生8个最小项：Y0=m0, Y1=m1, ..., Y7=m7
  \item $F = m1+m3+m5+m7 = Y1+Y3+Y5+Y7$
  \item 用1个4输入OR门连接Y1,Y3,Y5,Y7
\end{enumerate}

\textbf{【电路结构】}
\begin{center}
\begin{tikzpicture}
\node[draw,minimum width=2.5cm,minimum height=1.5cm] (dec) at (0,0) {3-8译码器};
\node[american or port] (orgate) at (4,0) {};
\draw[->] (dec.east|-orgate.in 1) -- (orgate.in 1);
\draw[->] (dec.east|-orgate.in 2) -- (orgate.in 2);
\draw[->] (orgate.out) -- ++(1,0) node[right]{$F$};
\node at (0,-2) {A2A1A0 = 变量A,B,C};
\node at (4,-2) {OR门输入=Y1,Y3,Y5,Y7};
\end{tikzpicture}
\end{center}

\textbf{【VHDL】}
\begin{lstlisting}
-- 例化译码器
dec: decoder_3to8 port map(A=>ABC, Y=>Y);
-- OR门
F <= Y(1) or Y(3) or Y(5) or Y(7);
\end{lstlisting}

\textbf{【推广：实现任意3变量函数】}
\begin{center}
\begin{tabular}{|l|l|}
\hline
\textbf{函数} & \textbf{连接} \\
\hline
$F=m0+m2+m4+m6$ & Y0+Y2+Y4+Y6 \\
$F=m0+m1+m2+m3$ & Y0+Y1+Y2+Y3 \\
$F=A$ (即m4+m5+m6+m7) & Y4+Y5+Y6+Y7 \\
$F=\overline{A}$ & Y0+Y1+Y2+Y3 \\
\hline
\end{tabular}
\end{center}

\textbf{【如果函数用NAND实现】}
用NAND门替代OR门（德摩根定理）。此时输出是反相的，但功能等价。

\begin{examfocus}
\textbf{必考！用3-8译码器+门电路实现任意3变量逻辑函数。}
步骤：(1)将函数写成最小项之和 (2)对应Y输出用OR门连接
\end{examfocus}

\end{detailbox}


\section{变式3：用3-8译码器构建4-16译码器（片选级联）}

\begin{detailbox}[完整教学]

\textbf{【问题】}只有3-8译码器，如何得到4位地址→16个输出的4-16译码器？

\textbf{【方案】}两片3-8译码器+最高位地址A3做片选。

\textbf{【级联逻辑】}
\begin{itemize}
  \item \textbf{片0（低位）}：处理A3=0时，地址0000-0111（译码器输出Y0-Y7对应系统Y0-Y7）
  \item \textbf{片1（高位）}：处理A3=1时，地址1000-1111（译码器输出Y0-Y7对应系统Y8-Y15）
  \item A3=0→片0使能，片1禁止
  \item A3=1→片1使能，片0禁止
\end{itemize}

\textbf{【VHDL】}
\begin{lstlisting}
entity decoder_4to16 is
  port (A : in std_logic_vector(3 downto 0);
        Y : out std_logic_vector(15 downto 0));
end decoder_4to16;

architecture structural of decoder_4to16 is
  signal Y_low, Y_high : std_logic_vector(7 downto 0);
  component decoder_3to8 is
    port (A : in std_logic_vector(2 downto 0);
          E : in std_logic; Y : out std_logic_vector(7 downto 0));
  end component;
begin
  -- 片0：A3=0时使能
  dec_low: decoder_3to8 port map(
    A => A(2 downto 0), E => not A(3), Y => Y_low);

  -- 片1：A3=1时使能
  dec_high: decoder_3to8 port map(
    A => A(2 downto 0), E => A(3), Y => Y_high);

  Y <= Y_high & Y_low;  -- 拼接为16位输出
end structural;
\end{lstlisting}

\textbf{【通用级联公式】}
用$2^k$个$n$-$2^n$译码器构建$(n+k)$-$2^{n+k}$译码器。
高k位地址→片选，低n位地址→各片输入。

\textbf{【常见考题】}用两片74LS138构建4-16译码器。利用使能端做片选。

\end{detailbox}


\section{变式4：用译码器做地址译码}

\begin{detailbox}[完整教学]

\textbf{【应用场景】}
CPU有16位地址总线，但连接了多个外设（RAM、ROM、IO端口等）。
用译码器将地址空间划分为不同区域，产生各外设的片选信号。

\textbf{【例题】}用3-8译码器产生8个外设的片选信号。地址A15-A13做译码输入。

\textbf{【电路】}
\begin{center}
\begin{tabular}{|c|c|l|}
\hline
A15A14A13 & 选中输出 & 外设 \\
\hline
000 & Y0 & RAM1(0x0000-0x1FFF) \\
001 & Y1 & RAM2(0x2000-0x3FFF) \\
010 & Y2 & ROM (0x4000-0x5FFF) \\
011 & Y3 & IO1 (0x6000-0x7FFF) \\
100 & Y4 & IO2 (0x8000-0x9FFF) \\
\hline
\end{tabular}
\end{center}

\textbf{【核心思想】}高位地址→译码→片选。低位地址→片内寻址。

\end{detailbox}


\section{变式5：显示译码器（BCD→7段）}

\begin{detailbox}[完整教学]

\textbf{【问题】}3-8译码器每次只选中1个输出。但7段数码管需要同时点亮多个段（如数字"8"需要7段全亮）。

\textbf{【区别】}
\begin{itemize}
  \item 3-8译码器：$n$输入→$2^n$输出，输出是\textbf{互斥}的（恰好一个为1）
  \item BCD-7段译码器：4输入→7输出，输出\textbf{不互斥}（可以同时多个为1）
\end{itemize}

\textbf{【详见第三章BCD译码器】}BCD-7段译码器本质上是7个独立的组合逻辑函数，不是简单的"最小项展开"。

\begin{examfocus}
\textbf{关键区别}：
\begin{itemize}
  \item 3-8译码器：1-of-N选择（互斥输出）→ 用AND门
  \item BCD-7段：任意多段同时亮（非互斥）→ 用ROM/查找表
\end{itemize}
不搞混这两种"译码器"是考试的基本要求。
\end{examfocus}

\end{detailbox}


\section{译码器变式总结}

\begin{examfocus}[译码器5大变式速查]
\begin{center}
\begin{tabular}{|l|l|l|}
\hline
\textbf{变式} & \textbf{核心考点} & \textbf{常考题型} \\
\hline
74LS138风格 & 低有效+3使能端 & 写VHDL, 判断使能条件 \\
译码器+OR≈任意函数 & 最小项之和 & 给定函数→画电路 \\
级联构建更大译码器 & 高位地址做片选 & 74LS138×2→4-16 \\
地址译码 & 地址空间划分 & 存储器片选设计 \\
vs BCD-7段 & 互斥vs非互斥输出 & 概念辨析 \\
\hline
\end{tabular}
\end{center}
\end{examfocus}

\chapter{BCD译码器——5大变式详解}

\section{变式1：多位BCD动态扫描显示}

\begin{detailbox}[完整教学]
\textbf{【问题】}4位数码管需要4个BCD-7段译码器吗？如果每个译码器单独驱动→需要4×7=28个IO口。

\textbf{【方案】}动态扫描：1个译码器驱动所有数码管的段，快速轮流点亮各数码管（>50Hz），利用人眼视觉暂留。

\textbf{【电路】}
计数器→2-4译码器→轮流选通4个数码管的共阴/共阳极。BCD译码器输出段码到所有数码管。

\textbf{【VHDL核心逻辑】}
\begin{lstlisting}
process(clk)  -- 扫描时钟(通常>200Hz)
begin
  if rising_edge(clk) then
    scan_cnt <= scan_cnt + 1;
    case scan_cnt is
      when "00" => digit_sel <= "0001"; seg_data <= digit0;
      when "01" => digit_sel <= "0010"; seg_data <= digit1;
      when "10" => digit_sel <= "0100"; seg_data <= digit2;
      when "11" => digit_sel <= "1000"; seg_data <= digit3;
    end case;
  end if;
end process;
\end{lstlisting}

\textbf{【常见考题】}设计4位动态扫描显示电路。
\end{detailbox}


\section{变式2：带锁存的BCD译码器}

\begin{detailbox}[完整教学]
\textbf{【问题】}BCD译码器输入随时钟快速变化→数码管闪烁。需要"定格"显示。

\textbf{【方案】}在译码器前加寄存器。LE=1时锁存当前BCD值，LE=0时保持上次锁存值不变。

\textbf{【VHDL】}
\begin{lstlisting}
process(clk)
begin
  if rising_edge(clk) then
    if LE = '1' then
      latched_bcd <= bcd_input;  -- 锁存
    end if;
  end if;
end process;
seg <= bcd_to_7seg(latched_bcd);  -- 译码锁存后的值
\end{lstlisting}

\textbf{【设计思想】}锁存=寄存器=1个时钟沿保存数据。本质=在数据通路上插入一级寄存器。
\end{detailbox}


\section{变式3：共阳极 vs 共阴极数码管}

\begin{detailbox}[完整教学]
\textbf{【共阳极】}所有LED阳极接VCC。段驱动需低电平(0)点亮LED。seg='0'=亮。
\textbf{【共阴极】}所有LED阴极接GND。段驱动需高电平(1)点亮LED。seg='1'=亮。

\textbf{【输出值关系】}$seg_{\text{共阳}} = \text{not } seg_{\text{共阴}}$（互为反码）

\textbf{【VHDL适配】}
\begin{lstlisting}
-- 共阴极(1亮)
when "0000" => seg <= "1111110";  -- 0
-- 共阳极(0亮)
when "0000" => seg <= "0000001";  -- 0(反码)
\end{lstlisting}

\textbf{【考试陷阱】}题目未明确说明共阴/共阳时，看段码值判断：1111110(共阴0) vs 0000001(共阳0)。
\end{detailbox}


\section{变式4：计数器+译码器+数码管——综合题}

\begin{detailbox}[完整教学]
\textbf{【经典考题】}设计0-9计数器+BCD-7段译码器+数码管显示。

\textbf{【完整VHDL】}
\begin{lstlisting}
entity counter_display is
  port (clk, rst_n : in std_logic;
        seg : out std_logic_vector(6 downto 0));
end counter_display;

architecture rtl of counter_display is
  signal cnt : std_logic_vector(3 downto 0);
  function bcd7seg(bcd : std_logic_vector(3 downto 0))
    return std_logic_vector is
  begin
    case bcd is
      when "0000" => return "1111110";
      when "0001" => return "0110000";
      when "0010" => return "1101101";
      when "0011" => return "1111001";
      when "0100" => return "0110011";
      when "0101" => return "1011011";
      when "0110" => return "1011111";
      when "0111" => return "1110000";
      when "1000" => return "1111111";
      when "1001" => return "1111011";
      when others => return "0000000";
    end case;
  end function;
begin
  process(clk, rst_n)
  begin
    if rst_n='0' then cnt <= "0000";
    elsif rising_edge(clk) then
      if cnt = "1001" then cnt <= "0000";
      else cnt <= cnt + 1;
      end if;
    end if;
  end process;
  seg <= bcd7seg(cnt);
end rtl;
\end{lstlisting}

\textbf{【考点】}两个模块整合：计数器+译码器。
\end{detailbox}


\section{变式5：16进制显示译码器（0-F）}

\begin{detailbox}[完整教学]
\textbf{【问题】}BCD译码器只显示0-9。如何显示0-F（全4位二进制）？

\textbf{【方案】}扩展case覆盖1010-1111，定义A-F的段码。

\begin{lstlisting}
-- 扩展BCD译码器支持0-F
when "1010" => return "1110111";  -- A
when "1011" => return "0011111";  -- b
when "1100" => return "1001110";  -- C
when "1101" => return "0111101";  -- d
when "1110" => return "1001111";  -- E
when "1111" => return "1000111";  -- F
\end{lstlisting}

\textbf{【考试】}有时问A-F的段码值，需要记住典型表示法。
\end{detailbox}

\chapter{4选1 MUX——4大变式详解}

\section{变式1：用MUX实现任意逻辑函数（必考！）}

\begin{detailbox}[完整教学]
\textbf{【核心原理】}
MUX的通用表达式：$Y=\overline{S_1}\cdot\overline{S_0}\cdot D_0+\overline{S_1}\cdot S_0\cdot D_1+S_1\cdot\overline{S_0}\cdot D_2+S_1\cdot S_0\cdot D_3$

如果把$S_1,S_0$接函数变量，$D_k$接"该最小项是否参与=1"：$D_k$接1→该最小项参与函数。$D_k$接0→不参与。

\textbf{【结论】任何n变量逻辑函数=一个$2^n$选1 MUX}

\textbf{【例题1】}用8选1 MUX实现$F(A,B,C)=\sum m(0,3,5,6)$

\textbf{解：}$D_0=1,D_1=0,D_2=0,D_3=1,D_4=0,D_5=1,D_6=1,D_7=0$

将m0-m7对应的D分别接1或0。选择线S2,S1,S0接A,B,C。

\textbf{【例题2：降维法】}用4选1 MUX实现3变量函数$F(A,B,C)=\sum m(0,1,4,6,7)$

$2^{3-1}=4$选1 MUX, 变量A,B接S1,S0, C用于产生$D_k$。

\begin{center}
\begin{tabular}{|c|c|c|c|c|}
\hline
AB & m(C=0) & m(C=1) & D输入 & 逻辑 \\
\hline
00 & m0=1 & m1=1 & D0 & 1(与C无关) \\
01 & m2=0 & m3=0 & D1 & 0 \\
10 & m4=1 & m5=0 & D2 & $\overline{C}$ \\
11 & m6=1 & m7=1 & D3 & 1 \\
\hline
\end{tabular}
\end{center}

D0=1, D1=0, D2=$\overline{C}$, D3=1.

\textbf{【VHDL】}
\begin{lstlisting}
entity mux_logic is
  port (A, B, C : in std_logic; F : out std_logic);
end mux_logic;
architecture rtl of mux_logic is
  signal D : std_logic_vector(3 downto 0);
  signal S : std_logic_vector(1 downto 0);
begin
  S <= A & B;
  D <= '1' & (not C) & '0' & '1';  -- D3,D2,D1,D0
  with S select F <=
    D(0) when "00", D(1) when "01",
    D(2) when "10", D(3) when "11";
end rtl;
\end{lstlisting}

\begin{examfocus}
\textbf{这是组合逻辑的必考题！}两种考法：
(1) $2^n$选1 MUX实现n变量函数 → Dk直接接0/1
(2) $2^{n-1}$选1 MUX实现n变量函数 → 需要降维，剩余变量或其反接到D端
\end{examfocus}
\end{detailbox}


\section{变式2：双4选1 MUX构建8选1 MUX}

\begin{detailbox}[完整教学]
\textbf{【结构】}两片4选1 MUX处理低2位地址。一片2选1 MUX用最高位地址S2选择输出。
总MUX数：3片(2×4选1 + 1×2选1)。

\textbf{【VHDL——结构级】}
\begin{lstlisting}
-- 片0：S1S0选D0-D3
mux_low: entity work.mux4to1 port map(D=>D(3 downto 0), S=>S(1 downto 0), Y=>F_low);
-- 片1：S1S0选D4-D7
mux_high: entity work.mux4to1 port map(D=>D(7 downto 4), S=>S(1 downto 0), Y=>F_high);
-- 顶层2选1
F <= F_low when S(2)='0' else F_high;
\end{lstlisting}
\end{detailbox}


\section{变式3：带使能端的MUX}

\begin{detailbox}[完整教学]
\textbf{【功能】}E=1正常选择；E=0输出高阻态'Z'（三态输出）或固定'0'。

\textbf{【VHDL——三态输出】}
\begin{lstlisting}
F <= D(to_integer(unsigned(S))) when E='1' else 'Z';
\end{lstlisting}

\textbf{【应用】}多个MUX输出可共用一根总线（总线仲裁）。
\end{detailbox}


\section{变式4：MUX扩展——16选1}

\begin{detailbox}[完整教学]
\textbf{【树形结构】}4片4选1 MUX(第一层)+1片4选1 MUX(第二层)=16选1。
总MUX数：5片。

\textbf{【选择位分配】}S1S0→第一层(4片)，S3S2→第二层(1片)。
\end{detailbox}

\chapter{4位比较器——4大变式详解}

\section{变式1：8位比较器（两片4位级联）}

\begin{detailbox}[完整教学]
\textbf{【级联方式】}
高位4位→主比较器。低位4位→从比较器，输出接主比较器的级联输入端。

\textbf{【级联逻辑】}
\begin{itemize}
  \item 主比较器比较A[7:4]和B[7:4]
  \item 如果高位相等，则看级联输入（低位比较器的结果）
  \item 低位比较器比较A[3:0]和B[3:0]，输出接到主比较器的I\_A>B, I\_A=B, I\_A<B
\end{itemize}

\textbf{【VHDL——结构级】}
\begin{lstlisting}
-- 低位比较器
cmp_lo: entity work.comparator_4bit
  port map(A=>A(3 downto 0), B=>B(3 downto 0),
           I_gt=>'0', I_eq=>'1', I_lt=>'0',  -- 最低位：假设之前相等
           A_gt=>gt_lo, A_eq=>eq_lo, A_lt=>lt_lo);
-- 高位比较器
cmp_hi: entity work.comparator_4bit
  port map(A=>A(7 downto 4), B=>B(7 downto 4),
           I_gt=>gt_lo, I_eq=>eq_lo, I_lt=>lt_lo,  -- 级联!
           A_gt=>A_gt, A_eq=>A_eq, A_lt=>A_lt);
\end{lstlisting}

\textbf{【行为级——直接用运算符（更简洁）】}
\begin{lstlisting}
A_gt <= '1' when A > B else '0';
A_eq <= '1' when A = B else '0';
A_lt <= '1' when A < B else '0';
\end{lstlisting}
综合器自动处理任意位宽的比较！
\end{detailbox}


\section{变式2：只输出"相等"的比较器}

\begin{detailbox}[完整教学]
\textbf{【简化】}只需要$A=B$信号。每一位都相等→XNOR→全部AND。

\textbf{【VHDL】}
\begin{lstlisting}
A_eq <= '1' when A = B else '0';
-- 或门级：
A_eq <= (A(3) xnor B(3)) and (A(2) xnor B(2)) and
        (A(1) xnor B(1)) and (A(0) xnor B(0));
\end{lstlisting}

\textbf{【应用】}地址译码、标签匹配、Cache Tag比较。
\end{detailbox}


\section{变式3：比较器+MUX实现取最大值}

\begin{detailbox}[完整教学]
\textbf{【电路】}比较器判断A>B→输出作为MUX选择信号。A>B时MUX选A，否则选B。

\textbf{【VHDL】}
\begin{lstlisting}
max_val <= A when A > B else B;
min_val <= A when A < B else B;
\end{lstlisting}

\textbf{【综合结果】}比较器+MUX→单行代码综合成两者组合。
\end{detailbox}


\section{变式4：带使能的比较器}

\begin{detailbox}[完整教学]
\textbf{【功能】}enable=1时比较器工作；enable=0时所有输出=0。

\textbf{【VHDL】}
\begin{lstlisting}
process(A, B, en)
begin
  if en = '0' then
    A_gt <= '0'; A_eq <= '0'; A_lt <= '0';
  else
    if A > B then A_gt <= '1'; A_eq <= '0'; A_lt <= '0';
    elsif A = B then A_gt <= '0'; A_eq <= '1'; A_lt <= '0';
    else A_gt <= '0'; A_eq <= '0'; A_lt <= '1';
    end if;
  end if;
end process;
\end{lstlisting}
\end{detailbox}

\chapter{4位加法器——5大变式详解}

\section{变式1：8位/16位加法器（级联）}

\begin{detailbox}[完整教学]
\textbf{【方案】}两片4位加法器级联。低位$C_{out}$→高位$C_{in}$。

\textbf{【VHDL——结构级】}
\begin{lstlisting}
add_lo: entity work.adder_4bit port map(A=>A(3 downto 0), B=>B(3 downto 0),
        Cin=>'0', Sum=>Sum(3 downto 0), Cout=>carry);
add_hi: entity work.adder_4bit port map(A=>A(7 downto 4), B=>B(7 downto 4),
        Cin=>carry, Sum=>Sum(7 downto 4), Cout=>Cout);
\end{lstlisting}

\textbf{【行为级——推荐】}
\begin{lstlisting}
signal result : std_logic_vector(8 downto 0);
result <= ('0' & A) + ('0' & B);
Sum <= result(7 downto 0); Cout <= result(8);
\end{lstlisting}
综合器自动优化任意位宽加法。
\end{detailbox}


\section{变式2：减法器（用加法器实现A-B）}

\begin{detailbox}[完整教学]
\textbf{【补码原理】}$A-B = A+(-B) = A+(\sim B+1) = A+\sim B+1$

\textbf{【电路】}B取反→接加法器B输入。$C_{in}=1$（实现+1）。$C_{in}$就是补码中的"+1"。

\textbf{【VHDL】}
\begin{lstlisting}
diff <= A + (not B) + '1';  -- A-B = A + (~B) + 1
-- 或直接用减法运算符
diff <= A - B;
\end{lstlisting}

\begin{examfocus}
\textbf{常考}：为什么$C_{in}=1$能实现减法？
答案：$A-B = A+(\sim B)+1$ → +1部分由$C_{in}=1$提供。
\end{examfocus}
\end{detailbox}


\section{变式3：加法器+寄存器=累加器}

\begin{detailbox}[完整教学]
\textbf{【结构】}加法器计算$ACC+Input$。结果在时钟沿存入ACC。ACC输出反馈回加法器输入端。

\textbf{【VHDL】}
\begin{lstlisting}
process(clk, rst_n)
begin
  if rst_n='0' then acc <= (others=>'0');
  elsif rising_edge(clk) then
    acc <= acc + input_data;  -- 每个时钟累加一次
  end if;
end process;
\end{lstlisting}

\textbf{【本质】}累加器=计数器的一种（不是每次+1，而是+任意值）。

\textbf{【应用】}DSP中的乘累加(MAC)、积分器。
\end{detailbox}


\section{变式4：带溢出检测的加法器}

\begin{detailbox}[完整教学]
\textbf{【无符号数溢出】}$C_{out}=1$=溢出（结果超过最大值）。

\textbf{【有符号数溢出】}$Overflow = C_n \oplus C_{n-1}$（最高位进位和次高位进位异或）。

\textbf{【VHDL】}
\begin{lstlisting}
signal result : std_logic_vector(4 downto 0);  -- 5位
result <= ('0' & A) + ('0' & B);
Sum <= result(3 downto 0);
Cout <= result(4);
-- 有符号溢出
Overflow <= result(4) xor result(3);  -- C4 xor C3
\end{lstlisting}

\textbf{【直观判断】}
两个正数相加得负数 → 溢出 (正溢出)
两个负数相加得正数 → 溢出 (负溢出)
\end{detailbox}


\section{变式5：BCD加法器}

\begin{detailbox}[完整教学]
\textbf{【问题】}二进制加法后，如果结果>9(1001)，需要+6(0110)修正为BCD格式。

\textbf{【VHDL】}
\begin{lstlisting}
signal bin_sum : std_logic_vector(4 downto 0);
signal bcd_sum : std_logic_vector(3 downto 0);
begin
  bin_sum <= ('0' & A) + ('0' & B);
  bcd_sum <= bin_sum(3 downto 0) when bin_sum <= 9
        else bin_sum(3 downto 0) + "0110";  -- >9时+6修正
  Cout <= '1' when bin_sum > 9 else '0';
\end{lstlisting}

\textbf{【例子】}5+7=12(1100)。1100+0110=10010→BCD: 0001 0010(12)。
\end{detailbox}

\chapter{级联MUX——4大变式详解}

\section{变式1：用2选1 MUX构建8选1 MUX}

\begin{detailbox}[完整教学]
\textbf{【树形结构】}3层：第一层4个→第二层2个→第三层1个=7个2选1 MUX。

\textbf{【选择位分配】}
\begin{itemize}
  \item S0→第一层(4个MUX)：每两组数据中选一个
  \item S1→第二层(2个MUX)：第一层结果中两两选择
  \item S2→第三层(1个MUX)：第二层两个结果中最终选择
\end{itemize}

\textbf{【通用公式】}$2^n$选1 MUX需要$2^n-1$个2选1 MUX，分$n$层。

\textbf{【VHDL——generate结构级】}
\begin{lstlisting}
-- 递归描述流水线MUX树
signal layer0 : std_logic_vector(3 downto 0);
signal layer1 : std_logic_vector(1 downto 0);
begin
  gen0: for i in 0 to 3 generate
    layer0(i) <= D(2*i) when S(0)='0' else D(2*i+1);
  end generate;
  gen1: for i in 0 to 1 generate
    layer1(i) <= layer0(2*i) when S(1)='0' else layer0(2*i+1);
  end generate;
  Y <= layer1(0) when S(2)='0' else layer1(1);
end behavioral;
\end{lstlisting}
\end{detailbox}


\section{变式2：用4选1 MUX构建16选1 MUX}

\begin{detailbox}[完整教学]
\textbf{【结构】}第一层4片4选1 MUX，第二层1片4选1 MUX。共5片。

S1S0→第一层(4片4选1 MUX)。S3S2→第二层(1片4选1 MUX选择第一层输出之一)。
\end{detailbox}


\section{变式3：不对称级联（构建6选1 MUX）}

\begin{detailbox}[完整教学]
\textbf{【问题】}需要6选1，但MUX只有$2^n$选1。用8选1实现6选1（两个输入悬空或接地）。

\textbf{【VHDL——4选1+2选1】}
\begin{lstlisting}
-- 第一层：4选1处理D0-D3
mux4: entity work.mux4to1 port map(D=>D(3 downto 0), S=>S(1 downto 0), Y=>mux4_out);
-- 顶层：2选1在mux4_out和D4/D5间选择
Y <= mux4_out when S(2)='0' else
     D(4) when S(1 downto 0)="00" else D(5);
-- 注意：S2=1时S1S0只需表示2种状态即可
\end{lstlisting}

\textbf{【简化方案】}直接用8选1 MUX，D6和D7接地(固定为0)。
\end{detailbox}


\section{变式4：MUX树实现逻辑函数}

\begin{detailbox}[完整教学]
\textbf{【本质】}MUX=通用逻辑实现器。MUX树=多层选择=分而治之。

\textbf{【Shannon展开定理】}$F(x_1,x_2,\dots,x_n)= \overline{x_1}\cdot F(0,x_2,\dots,x_n) + x_1\cdot F(1,x_2,\dots,x_n)$

这正是2选1 MUX的表达式！递归应用=构建MUX树。

\textbf{【结论】}任何组合逻辑函数都可用2选1 MUX树实现——这其实就是FPGA中LUT的工作原理。
\end{detailbox}


% ============================================================
\part{D触发器变式详解}
\chapter{D触发器——核心变式详解}

\section{变式1：带异步复位的D触发器 vs 带同步复位的D触发器}

\begin{detailbox}[完整教学]
\textbf{【异步复位】}复位不受时钟控制，rst=0立即清零（不等时钟沿）。
\textbf{【同步复位】}复位只在时钟沿有效，rst=0后需等下一个时钟上升沿才清零。

\textbf{【VHDL对比】}
\begin{lstlisting}
-- 异步复位（rst在敏感列表中，且在rising_edge之前判断）
process(clk, rst_n)
begin
  if rst_n = '0' then Q <= '0';           -- 立即响应
  elsif rising_edge(clk) then Q <= D;
  end if;
end process;

-- 同步复位（rst不在敏感列表，在rising_edge内部判断）
process(clk)
begin
  if rising_edge(clk) then
    if rst_n = '0' then Q <= '0';         -- 等时钟沿
    else Q <= D;
    end if;
  end if;
end process;
\end{lstlisting}

\textbf{【区别】}
\begin{center}
\begin{tabular}{|l|c|c|}
\hline & 异步复位 & 同步复位 \\
\hline
敏感列表 & clk, rst\_n & clk \\
复位延迟 & 立即 & 等1个时钟周期 \\
抗噪性 & 差(毛刺可触发复位) & 好(只在沿采样) \\
综合结果 & DFF+异步CLR & DFF+同步CLR \\
\hline
\end{tabular}
\end{center}

\begin{examfocus}
\textbf{必考判断}：看process敏感列表有没有rst\_n。
\begin{itemize}
  \item 有→异步复位
  \item 没有→同步复位
\end{itemize}
\end{examfocus}
\end{detailbox}


\section{变式2：带时钟使能的D触发器}

\begin{detailbox}[完整教学]
\textbf{【功能】}CE=1时正常采样D；CE=0时保持Q不变。

\textbf{【VHDL】}
\begin{lstlisting}
process(clk)
begin
  if rising_edge(clk) then
    if CE = '1' then Q <= D; end if;
    -- CE='0'：不赋值=保持
  end if;
end process;
\end{lstlisting}

\textbf{【综合结果】}D触发器+时钟使能逻辑（MUX反馈Q或采样D）。

\textbf{【应用】}所有带使能的寄存器/计数器都基于这个结构。
\end{detailbox}


\section{变式3：多位D触发器=寄存器}

\begin{detailbox}[完整教学]
\textbf{【原理】}N个D触发器并行→N位寄存器。同一时钟沿采样N位输入。

\textbf{【VHDL】}
\begin{lstlisting}
signal reg : std_logic_vector(7 downto 0);
process(clk) begin
  if rising_edge(clk) then reg <= data_in; end if;
end process;
\end{lstlisting}

\textbf{【综合】}8个D触发器→8个FPGA寄存器。

\textbf{【应用】}流水线寄存器、数据缓冲、状态寄存器。
\end{detailbox}


\section{变式4：D触发器构成2分频器}

\begin{detailbox}[完整教学]
\textbf{【原理】}$\overline{Q}$接回$D$→$Q(n+1)=\overline{Q(n)}$→每拍翻转→2分频。

\textbf{【VHDL】}
\begin{lstlisting}
process(clk) begin
  if rising_edge(clk) then Q <= not Q; end if;
end process;
\end{lstlisting}

\textbf{【这是分频器和计数器的原子单元。】}
\end{detailbox}


\section{变式5：建立/保持时间与最大工作频率}

\begin{detailbox}[完整教学]
\textbf{【建立时间$t_{su}$】}时钟沿到来前，D输入必须稳定多久。
\textbf{【保持时间$t_h$】}时钟沿到来后，D输入必须保持多久。
\textbf{【$t_{cko}$】}时钟沿到Q输出有效的延迟。

\textbf{【最大频率】}$f_{max} = 1/(t_{cko} + t_{comb} + t_{su})$

$t_{comb}$=触发器之间的组合逻辑延迟。

\textbf{【考试计算】}给定$t_{cko}$=2ns, $t_{su}$=1ns, 组合逻辑延迟=5ns。
$f_{max}=1/(2+5+1)=125$MHz。
\end{detailbox}


% ============================================================
\part{计数器变式详解}
\chapter{计数器——统一框架变式详解}

\section{变式1：带使能的模N计数器}

\begin{detailbox}[完整教学]
\textbf{【电路】}在标准计数器上增加en控制。
\textbf{【VHDL】}
\begin{lstlisting}
if rising_edge(clk) then
  if en = '1' then
    if cnt = N-1 then cnt <= 0; else cnt <= cnt + 1; end if;
  end if;
end if;
\end{lstlisting}
\textbf{【考点】}en=0时不赋值→保持。这是所有可控计数器的基本范式。
\end{detailbox}


\section{变式2：带同步加载的模N计数器}

\begin{detailbox}[完整教学]
\textbf{【功能】}load=1时cnt=data\_in。load=0时正常计数。
\textbf{【优先级】}load > 计数（if-else中load在前=优先级高）。
\textbf{【VHDL】}
\begin{lstlisting}
if load = '1' then cnt <= data_in;
elsif cnt = N-1 then cnt <= 0;
else cnt <= cnt + 1; end if;
\end{lstlisting}
\end{detailbox}


\section{变式3：递增/递减双向计数器}

\begin{detailbox}[完整教学]
\textbf{【VHDL】}
\begin{lstlisting}
if up = '1' then cnt <= cnt + 1; else cnt <= cnt - 1; end if;
\end{lstlisting}
\textbf{【递减自然溢出】}$0000-1=1111$（4位无符号），自动模16循环。
\end{detailbox}


\section{变式4：倒计时计数器（预置→0→预置）}

\begin{detailbox}[完整教学]
\textbf{【VHDL】}
\begin{lstlisting}
if cnt = 0 then cnt <= PRE;  -- 到0后重新加载
else cnt <= cnt - 1; end if;
\end{lstlisting}
\end{detailbox}


\section{变式5：级联计数器——总模值=乘积}

\begin{detailbox}[完整教学]
\textbf{【级联方式】}低位carry→高位en。总M=$M_1 \times M_2$。

\textbf{【实例】}模8+模5级联→总模40。模60(6×10)BCD级联。

\textbf{【VHDL关键】}高位以低位carry为时钟使能——确保高位只在低位满一轮时才+1。
\end{detailbox}


\section{变式6：格雷码计数器}

\begin{detailbox}[完整教学]
\textbf{【特点】}相邻状态仅1位不同。适合FIFO指针、跨时钟域。
\textbf{【VHDL】}case枚举所有格雷码状态。
\textbf{【3位格雷码序列】}000-001-011-010-110-111-101-100-000。
\end{detailbox}


% ============================================================
\part{分频器变式详解}
\chapter{分频器——全部变式详解}

\section{变式1：直接取Qn实现$2^n$分频}

\begin{detailbox}[完整教学]
\textbf{【原理】}$Q_k$每$2^{k+1}$拍完成一次0→1→0循环。

\textbf{【VHDL——一行搞定】}
\begin{lstlisting}
clk_div2 <= cnt(0);   -- 2分频
clk_div4 <= cnt(1);   -- 4分频
clk_div8 <= cnt(2);   -- 8分频
\end{lstlisting}

\textbf{【优点】}不需要任何额外逻辑。计数器本身就是分频器。
\end{detailbox}


\section{变式2：偶数分频（任意M）}

\begin{detailbox}[完整教学]
\textbf{【统一模板】}
\begin{lstlisting}
if cnt = M-1 then cnt <= 0; else cnt <= cnt + 1; end if;
clk_out <= '1' when cnt >= M/2 else '0';  -- 50\%占空比
\end{lstlisting}
\end{detailbox}


\section{变式3：奇数分频（双沿法）}

\begin{detailbox}[完整教学]
\textbf{【原理】}上升沿+下降沿各产生一个信号→OR组合。得到接近50\%占空比的奇数分频。
\textbf{【详见解析册题6】}
\end{detailbox}


\section{变式4：占空比可调分频器}

\begin{detailbox}[完整教学]
\textbf{【原理】}改变输出翻转的阈值。占空比=高拍数/M。
\begin{lstlisting}
clk_out <= '1' when cnt < DUTY else '0';
\end{lstlisting}
\end{detailbox}


\section{变式5：多级分频器（分频比相乘）}

\begin{detailbox}[完整教学]
\textbf{【原理】}M1分频→M2分频=总M1×M2分频。
\textbf{【注意】}后级以前级输出为时钟→跨时钟域风险。FPGA中推荐用使能信号而非门控时钟。
\end{detailbox}


\section{变式6：秒脉冲发生器}

\begin{detailbox}[完整教学]
\textbf{【问题】}50MHz→1Hz。M=50,000,000。需26位计数器。
\textbf{【VHDL】}同偶数分频模板，M=50000000。
\end{detailbox}


% ============================================================
\part{序列发生器变式详解}
\chapter{序列发生器——全部变式详解}

\section{变式1：任意周期序列（计数器+映射表）}

\begin{detailbox}[完整教学]
\textbf{【通用方案】}模L计数器+状态→输出映射表。
\textbf{【例题】}产生序列10110(长度5)。模5计数器+case映射：
\begin{lstlisting}
case cnt is
  when "000" => seq <= '1';
  when "001" => seq <= '0';
  when "010" => seq <= '1';
  when "011" => seq <= '1';
  when "100" => seq <= '0';
end case;
\end{lstlisting}
\textbf{【适用范围】}任意序列。长度L→模L计数器。
\end{detailbox}


\section{变式2：状态机型序列发生器}

\begin{detailbox}[完整教学]
\textbf{【方案】}每个状态输出一个序列值。状态转移=序列顺序。
\textbf{【优点】}状态名有意义（如"S\_1","S\_0","S\_1","S\_1","S\_0"）。但状态多时代码冗长。
\textbf{【比较】}推荐用计数器+case——代码更简洁。
\end{detailbox}


\section{变式3：移位寄存器+LFSR发生器}

\begin{detailbox}[完整教学]
\textbf{【LFSR】}线性反馈移位寄存器。反馈=选择位的XOR。
\textbf{【应用】}伪随机序列(PRBS)、CRC、扰码。
\textbf{【4位LFSR(X4+X3+1)】}反馈=Q3 xor Q2。产生15个伪随机状态。
\textbf{【注意】}全0状态非法——LFSR会锁死在0000。必须初始化为非0值。
\end{detailbox}


\section{变式4：ROM/查找表型序列发生器}

\begin{detailbox}[完整教学]
\textbf{【方案】}预存整个序列在ROM中。计数器=地址。
\textbf{【优点】}任意序列、可重配置。适合长序列。
\textbf{【VHDL】}
\begin{lstlisting}
type rom_type is array (0 to 15) of std_logic_vector(3 downto 0);
constant SEQ_ROM : rom_type := (x"1", x"0", x"1", x"1", x"0", ...);
seq <= SEQ_ROM(to_integer(unsigned(cnt)));
\end{lstlisting}
\end{detailbox}


% ============================================================
\part{移位寄存器变式详解}
\chapter{移位寄存器——全部变式详解}

\section{变式1：串并转换（SIPO）}

\begin{detailbox}[完整教学]
\textbf{【功能】}串行输入→4拍后并行读取。UART/SPI接收器的基础。
\textbf{【VHDL】}\texttt{reg <= si \& reg(N-1 downto 1); po <= reg;}
\end{detailbox}


\section{变式2：并串转换（PISO）}

\begin{detailbox}[完整教学]
\textbf{【功能】}load并行数据→逐拍串行输出。UART/SPI发送器的基础。
\textbf{【VHDL】}
\begin{lstlisting}
if load='1' then reg <= data_in;
else reg <= '0' & reg(N-1 downto 1); end if;
so <= reg(0);
\end{lstlisting}
\end{detailbox}


\section{变式3：双向移位}

\begin{detailbox}[完整教学]
\textbf{【VHDL】}
\begin{lstlisting}
if dir='1' then reg <= si_r & reg(N-1 downto 1);
else reg <= reg(N-2 downto 0) & si_l; end if;
\end{lstlisting}
\end{detailbox}


\section{变式4：循环移位}

\begin{detailbox}[完整教学]
\textbf{【功能】}数据不丢失——最低位移回最高位。
\textbf{【VHDL】}\texttt{reg <= reg(0) \& reg(N-1 downto 1);}（右循环）
\textbf{【本质】}这就是环形计数器的基础！
\end{detailbox}


\section{变式5：算术移位（保持符号位）}

\begin{detailbox}[完整教学]
\textbf{【逻辑右移】}高位补0 → \texttt{'0' \& reg(N-1 downto 1)}
\textbf{【算术右移】}高位补符号位 → \texttt{reg(N-1) \& reg(N-1 downto 1)}
\textbf{【应用】}有符号数÷2操作。
\end{detailbox}


% ============================================================
\part{环形计数器变式详解}
\chapter{环形计数器——全部变式详解}

\section{变式1：N位环形计数器}

\begin{detailbox}[完整教学]
\textbf{【通用模板】}
\begin{lstlisting}
process(clk, rst_n)
begin
  if rst_n='0' then reg <= (N-1=>'1', others=>'0');  -- One-hot初值
  elsif rising_edge(clk) then
    reg <= reg(0) & reg(N-1 downto 1);  -- 右移+反馈
  end if;
end process;
\end{lstlisting}
\textbf{【结论】}N个DFF→N个状态。所有状态中恰好一个1（One-hot编码）。
\end{detailbox}


\section{变式2：约翰逊计数器（2N个状态）}

\begin{detailbox}[完整教学]
\textbf{【反馈】}\texttt{reg <= (not reg(0)) \& reg(N-1 downto 1);}
\textbf{【4位约翰逊状态序列】}$0000\to1000\to1100\to1110\to1111\to0111\to0011\to0001\to0000$
\textbf{【结论】}N个DFF→2N个状态（比环形计数器多一倍）。
\end{detailbox}


\section{变式3：自启动环形计数器}

\begin{detailbox}[完整教学]
\textbf{【问题】}干扰可能使环形计数器进入非法状态。
\textbf{【方案】}检测非法状态（如全0或多个1），强制恢复到合法状态(1000)。
\textbf{【VHDL】}在移位前检查：\texttt{if count\_ones(reg) /= 1 then reg <= "1000";}
\end{detailbox}


\section{变式4：用环形计数器实现时序控制器}

\begin{detailbox}[完整教学]
\textbf{【应用】}多周期操作器的步骤控制。
\textbf{【5步控制器】}5位环形计数器→每拍恰好一个步骤信号为1→控制对应步骤。
\textbf{【优点】}无需额外的译码逻辑——直接读Q值。
\end{detailbox}


% ============================================================
\part{序列检测器变式详解}
\chapter{序列检测器——全部变式详解}

\section{变式1：检测1011（重叠 vs 非重叠）}

\begin{detailbox}[完整教学]
\textbf{【核心区别】}S101+输入1→检测到！
\begin{itemize}
  \item \textbf{非重叠}→回S0（重新开始）
  \item \textbf{重叠}→去S1（最后的"1"作为新序列的第一个"1"）
\end{itemize}
\textbf{【考试判断】}看检测到后的转移目标。S0=非重叠，S1(或其他非S0状态)=重叠。
\end{detailbox}


\section{变式2：检测1101}

\begin{detailbox}[完整教学]
\textbf{【状态】}S0→S1→S11→S110→(输入1=检测到)
\textbf{【重叠处理】}S110+1→检测到→去S11（"1101"中后两位"11"匹配"11"前缀）
\textbf{【关键】}S11+1→S11（"111"中后两个"11"为有效前缀）。
\end{detailbox}


\section{变式3：检测1111（连续4个1）}

\begin{detailbox}[完整教学]
\textbf{【状态】}S0→S1(1个1)→S11(2)→S111(3)→(第4个1→检测)
\textbf{【特点】}任意时刻输入0→回S0（连续中断）。
\textbf{【重叠】}S111+1→S111（"1111"后3个"111"为有效前缀）。
\end{detailbox}


\section{变式4：检测0101}

\begin{detailbox}[完整教学]
\textbf{【状态】}S0→S01(0)→S010(01)→S0101(检测)
\textbf{【注意】}S0101+0→S01（重叠：最后的"0"作为新序列的第一个"0"）。
\end{detailbox}


\section{变式5：用移位寄存器+比较器替代状态机}

\begin{detailbox}[完整教学]
\textbf{【方案】}移位寄存器存最近N位→与目标序列比较。
\textbf{【VHDL——两行搞定】}
\begin{lstlisting}
reg <= din & reg(3 downto 1);
detected <= '1' when reg = "1011" else '0';
\end{lstlisting}
\textbf{【优点】}代码极简。不需推导状态转移。
\textbf{【注意】}自动实现重叠检测。
\textbf{【考试适用】}除非题目明确要求"用状态机设计"，否则移位寄存器方案是最佳选择。
\end{detailbox}


% ============================================================
\part{VHDL综合变式详解}
\chapter{VHDL综合——全部变式详解}

\section{变式1：if-else的综合结果}

\begin{detailbox}[完整教学]
\textbf{【综合规则】}
\begin{itemize}
  \item if-else级联→\textbf{优先级MUX链}。上层if优先级最高。
  \item 带clk的if→产生\textbf{寄存器}(D触发器)。
  \item 不带clk的if→产生\textbf{组合逻辑}。
  \item 不完整的if(缺少else)→可能产生\textbf{锁存器(Latch)}——通常应避免。
\end{itemize}

\textbf{【考试判断】}
\begin{lstlisting}
-- 组合逻辑MUX
if sel='1' then y <= a; else y <= b; end if;

-- 时序逻辑（寄存器）
if rising_edge(clk) then
  if en='1' then q <= d; end if;  -- 带使能的D触发器
end if;
\end{lstlisting}
\end{detailbox}


\section{变式2：case的综合结果}

\begin{detailbox}[完整教学]
\textbf{【综合规则】}
\begin{itemize}
  \item case→\textbf{并行MUX}或\textbf{译码器}。所有分支并行判断，无优先级。
  \item case不完备且无when others→可能产生锁存器。
  \item case完备(有when others)→纯组合逻辑或状态机。
\end{itemize}

\textbf{【if-else vs case硬件对比】}
\begin{center}
\begin{tabular}{|l|l|l|}
\hline & if-else & case \\
\hline
综合结果 & 优先级编码器/级联MUX & 并行MUX/译码器 \\
优先级 & 有(从上到下递减) & 无(并行) \\
速度 & 慢(级联延迟) & 快(并行) \\
面积 & 小(级联结构) & 大(并行结构) \\
\hline
\end{tabular}
\end{center}
\end{detailbox}


\section{变式3：signal vs variable}

\begin{detailbox}[完整教学]
\textbf{【signal(<=)】}赋值在process\textbf{结束时}更新。适合跨进程通信和寄存器。
\textbf{【variable(:=)】}赋值\textbf{立即}更新。适合process内部中间计算。

\textbf{【综合结果】}
\begin{itemize}
  \item signal在时序process中→D触发器
  \item signal在组合process中→连线
  \item variable→通常不产生额外硬件（被综合器优化为连线）
  \item 但variable在读取前先赋值→可能综合为寄存器
\end{itemize}

\textbf{【考试判断】}看赋值符号：<=（信号）vs :=（变量）。
\end{detailbox}


\section{变式4：process的综合结果}

\begin{detailbox}[完整教学]
\textbf{【process本身不产生硬件】}硬件由process\textbf{内部}代码决定。

\begin{center}
\begin{tabular}{|l|l|}
\hline
\textbf{process写法} & \textbf{综合结果} \\
\hline
process(clk)+rising\_edge & D触发器/寄存器 \\
process(all\_inputs)+无clk & 组合逻辑 \\
process(clk,rst)+异步rst & D触发器+异步复位 \\
不完整敏感列表 & 可能锁存器(应避免) \\
\hline
\end{tabular}
\end{center}
\end{detailbox}


\section{变式5：常见VHDL模式→硬件对照表}

\begin{detailbox}[完整教学]
\begin{center}
\begin{tabular}{|l|l|}
\hline
\textbf{VHDL模式} & \textbf{硬件} \\
\hline
\texttt{if rising\_edge(clk) then Q<=D;} & 1个D触发器 \\
\texttt{y <= a when sel='1' else b;} & 2选1 MUX \\
\texttt{with sel select y<=...} & 并行MUX \\
\texttt{case state is when...} & 状态机+译码器 \\
\texttt{reg <= si \& reg(N-1 downto 1);} & 移位寄存器 \\
\texttt{cnt <= cnt + 1;} & 加法器+寄存器 \\
\texttt{if cnt=N-1 then cnt<=0;} & 比较器(清零条件) \\
\texttt{for i in 0 to 7 generate} & 8个相同的硬件副本 \\
\hline
\end{tabular}
\end{center}

\begin{examfocus}
\textbf{VHDL综合考点总结：}
\begin{itemize}
  \item 有时钟边沿=有触发器
  \item if-else=优先级，case=并行
  \item signal<=延迟更新，variable:=立即更新
  \item process只是容器，内部代码决定硬件
  \item for-generate=复制硬件N份
\end{itemize}
\end{examfocus}
\end{detailbox}


\end{document}
